\documentclass[a4paper]{article}
\usepackage{exercise}
%um nur aufgaben zu zeigen
%\usepackage[noanswer]{exercise} 
\usepackage{../images/preamble}


\begin{document}
	\vspace*{-2cm}
	\parbox{4cm}{\vspace{0.4cm}\includegraphics[width=4cm]{../images/logo_scaled.pdf}}
	\parbox{10.6cm}{\setstretch{2.0} \centering{ \huge \textsf{Der ROLF-Aufgabenwettbewerb }}\\  \vspace*{-.5cm} }
\thispagestyle{empty}	
\vspace{1cm}
\noindent
\newline
\textsc{Willkommen} beim ROLF-Aufgabenwettbewerb!\\
Ab jetzt könnt ihr jeden Monat von uns neue Aufgaben erhalten, die euch helfen, mehr Physik zu verstehen!\\ 
Jede Aufgabe hat eine geheime Anzahl an Punkten, die wir erst festlegen, wenn wir die Aufgaben bewerten. Was aber gleich bleibt, ist ein festes, prozentuales Bewertungsschema, mit dem wir jede Lösung einschätzen. Das beinhaltet alle Aspekte, die wir für wichtig beim Aufgabenlösen halten. Der eine große Teil der Bewertung ist der (richtige) Ansatz zum Lösen der Aufgabe. Hier wird klar, ob ihr die Aufgabenstellung und das Problem verstanden habt, und ob ihr eine Idee habt, wie man die Aufgabe lösen kann. Der andere große Teil der Bewertung ist die Argumentation zur Lösung der Aufgabe. Hier sieht man, ob ihr mit euren Ansätzen auch Lösungen finden könnt. 
\newline
\begin{figure}[h]
	\centering
\begin{tabular}{l| l c}
	Ansatz: & & 40\%\\\hline
	& nachvollziehbar:&20\%\\
	&zielführend:&40\%\\
	&math. Formulierung :& 40\%\\\hline\hline
	Argumentation:& &60\% \\\hline
	&nachvollziehbar:&40\%\\
	&richtig:&60\%\\\hline
	
\end{tabular}
\end{figure}
Die Punktzahlen sollen dafür stehen, wie schwer die Aufgaben sind. Zusätzlich kann es sein, dass eine (falsche) Annahme das Problem sehr vereinfacht. Das wird sich in einem prozentuellen Abzug von der Gesamtpunktzahl pro Aufgabe niederschlagen.\\
An dieser Stelle ein Paar Tipps, die dafür sorgen, dass wir bessere Laune haben, wenn wir eure Aufgaben korrigieren:
\begin{itemize}
	\item Einheiten sind eure Freunde! - Wenn man eine Temperatur und einen Widerstand addiert, eine Fläche und eine Masse, ist das ein gutes Zeichen dafür, dass man etwas falsch macht.
	\item Termumformungen können lang und hässlich werden! - Schreibt am Besten \underline{alle} Gleichungen auf, mit denen ihr denkt, die Aufgabe lösen zu können, bevor ihr mit langen Termumformungen anfangt.
	\item Aufgabenstellungen sind nicht zum Spaß da! - Wenn ihr gegebene Größen oder Annahmen aus der Aufgabenstellung nicht braucht, seid ihr entweder wirklich schlau oder habt die Aufgabe nicht richtig gelöst/verstanden.
\end{itemize}
Wenn ihr die Aufgaben bearbeitet habt, könnt ihr die Bearbeitungen immer per Mail an \textbf{ \href{mailto:physikrolf@gmail.com}{physikrolf@gmail.com}} senden. Oder ihr gebt sie dem, der euch die Aufgaben gegeben hat. An den Tagen, an denen ihr die Aufgaben abgeben solltet (gerne auch früher, etwa 1x im Monat) bekommt ihr gleich die Musterlösungen für diese Serie und die nächsten Aufgaben. Sobald wir fertig mit dem Korrigieren sind, gibt es auch eine Rückmeldungen. Schickt bitte eure Bearbeitungen auch ein, wenn ihr denkt, eine Aufgabe nicht gelöst zu haben. Das hilft uns, in Zukunft bessere zu machen!\\
Alle Aufgabenserien mit Lösungen, die aktuelle Aufgabenserie und zusätzliches Material findet ihr auch auf unserer Website, \url{pankratius.github.io/rolf} .
Insgesamt hoffen wir, dass ihr mit den Aufgaben ein wenig Spaß habt, und vielleicht sogar etwas dazu lernt. Feedback dazu, wie gut die Aufgaben sind, könnt ihr gerne mit abgeben!
\end{document}